\documentclass[a4paper,12pt]{article}

% Pacotes essenciais
\usepackage[utf8]{inputenc}
\usepackage[T1]{fontenc}
\usepackage[brazilian]{babel}
\usepackage{amsmath, amssymb}
\usepackage{graphicx}
\usepackage{geometry}
\usepackage{hyperref}
\usepackage{float}
\usepackage{indentfirst}

% para por código com corzinha ui ui
\usepackage{minted}
\usepackage{xcolor}

\definecolor{bg}{rgb}{0.1,0.1,0.1}
\usemintedstyle{monokai}

\geometry{a4paper, margin=2.5cm}

\DeclareMathOperator{\sen}{sen}

\begin{document}

\begin{titlepage}
\centering

    UNIVERSIDADE FEDERAL DE SERGIPE\\
    CIDADE UNIVERSITÁRIA PROFESSOR JOSÉ ALOÍSIO DE CAMPOS

\vspace{4cm}

    FRANCISCO PASSOS DOS SANTOS ALVES

\vspace{4cm}

    FORMULAÇÃO MATEMÁTICA DE A DIGITAL CLOCKWORK
\vfill

    SÃO CRISTÓVÃO\\
    2026

\end{titlepage}

\clearpage
\pagenumbering{arabic}


\newpage
\tableofcontents
\newpage


\section{Introdução}

O \textit{digital-clockwork-simulator} é um relógio digital implementado inteiramente a partir de componentes lógicos discretos, como portas lógicas, flip-flops, comparadores e contadores. Este trabalho tem como objetivo transcrever, com rigor matemático, a lógica de funcionamento desses circuitos para um modelo formal, utilizando conceitos de linguagens formais para descrever o comportamento de cada componente e suas interações. Essa formalização estabelece uma ponte entre a implementação simulada em C++ e sua futura descrição em Verilog, possibilitando posteriormente sua verificação formal em Rocq.

A motivação deste trabalho é dupla. Primeiro, a construção de um modelo matemático explícito exige que decisões de projeto muitas vezes implícitas no hardware ou no código sejam definidas com precisão: o instante exato em que um comparador altera seu estado lógico, a convenção adotada para bordas de subida e descida de um sinal periódico, e a forma como os elementos sequenciais evoluem diante de um pulso de clock. O modelo formal representa então o circuito como um sistema bem definido, reduzindo as ambiguidades presentes na implementação.

Segundo, essa especificação é a base necessária para demonstrar propriedades do sistema. Antes de provar que uma implementação está correta, é preciso definir com precisão qual é o comportamento esperado do circuito e essa definição é exatamente a especificação construída neste documento, que descreve as operações, estados e propriedades que caracterizam o funcionamento dos componentes.

O documento segue a cadeia lógica de funcionamento do relógio digital: parte do sinal físico da rede elétrica de 60Hz, usado como referência temporal do sistema, e avança pelas etapas de aquisição, digitalização, divisão de frequência, contagem de pulsos e formação das unidades de tempo. Cada seção apresenta a modelagem matemática dos circuitos envolvidos e suas relações dentro do sistema, construindo gradualmente a especificação formal do relógio digital.


\newpage


\subsection{Aquisição e digitalização do sinal da rede elétrica}

A alimentação elétrica proveniente da rede é originalmente representada por um sinal analógico alternado, cuja tensão varia continuamente no tempo. Considerando uma rede elétrica residencial de tensão eficaz aproximada de \(125\)V e frequência de \(60\)Hz, o comportamento instantâneo da tensão pode ser modelado por uma função senoidal periódica:

\[
f : \mathbb{R} \to \mathbb{R}, \qquad f(t)=125\sen(120\pi t).
\]

O termo \(120\pi\) está relacionado à frequência angular da rede, obtida pela relação

\[
\omega = 2\pi f,
\]

onde \(f=60\)Hz representa a frequência de oscilação do sinal. Portanto,

\[
\omega = 2\pi(60)=120\pi \ \text{rad/s}.
\]

Dessa forma, a tensão completa um ciclo a cada

\[
T=\frac{1}{60}\text{s},
\]

satisfazendo a propriedade de periodicidade

\[
f\left(t+\frac{1}{60}\right)=f(t), \qquad \forall t\in\mathbb{R}.
\]

Embora esse sinal seja adequado para representar matematicamente a rede elétrica, circuitos digitais não trabalham diretamente com níveis contínuos de tensão. Portanto, é necessário realizar uma etapa de condicionamento do sinal, reduzindo sua amplitude para um nível compatível com a eletrônica digital.

Considerando uma redução ideal de tensão por um fator \(25\), a senoide de \(125\)V é convertida para uma senoide de amplitude \(5\)V, mantendo a mesma frequência e preservando a informação temporal do sinal original. Essa transformação pode ser descrita pela função:

\[
g:\mathbb{R}\to\mathbb{R}, \qquad
g(t)=\frac{f(t)}{25}=5\sen(120\pi t).
\]

A divisão por \(25\) altera apenas a amplitude do sinal, não modificando sua frequência nem seu período. Dessa maneira,

\[
g\left(t+\frac{1}{60}\right)=g(t), \qquad \forall t\in\mathbb{R}.
\]

O sinal \(g(t)\), apesar de possuir amplitude adequada para circuitos digitais, ainda é uma grandeza analógica, pois assume infinitos valores possíveis dentro do intervalo \([-5,5]\)V. Para que esse sinal possa ser utilizado como referência de clock em um sistema digital, é necessário convertê-lo para uma representação binária, contendo apenas dois estados possíveis.

Essa conversão é realizada por meio de um comparador de tensão ideal, cujo funcionamento consiste em verificar o sinal de entrada em relação a um valor de referência. Neste modelo, considera-se o limiar de decisão igual a \(0\)V. Assim, quando a senoide encontra-se no semiciclo positivo, o comparador produz nível lógico alto; quando encontra-se no semiciclo negativo ou exatamente no cruzamento por zero, produz nível lógico baixo.

A saída digital desse processo é representada pela função:

\[
q:\mathbb{R}\to\{0,1\},
\]

definida como:

\[
q(t)=
\begin{cases}
1, & \text{se } g(t)>0,\\
0, & \text{se } g(t)\leq0.
\end{cases}
\]

Essa função representa matematicamente a transformação de um sinal analógico periódico em um sinal digital periódico. Como o comparador apenas altera os níveis de amplitude e mantém os instantes de transição determinados pelos cruzamentos da senoide com o eixo de tempo, a frequência do sinal permanece igual à frequência da rede elétrica.

Consequentemente, \(q(t)\) possui também período

\[
T=\frac{1}{60},
\]

porém agora assumindo somente os valores discretos \(0\) e \(1\). Durante metade do período, correspondente ao semiciclo positivo da senoide, o sinal permanece em nível lógico alto; durante a outra metade, correspondente ao semiciclo negativo, permanece em nível lógico baixo.

Portanto, a função \(q(t)\) representa uma onda quadrada ideal de \(60\)Hz, obtida a partir da senoide da rede elétrica através de um processo de redução de amplitude e comparação por zero. Esse sinal digital pode então ser utilizado como referência de clock para sincronizar os elementos sequenciais do circuito.

\end{document}